% Source: https://artofproblemsolving.com/community/c3h1371507_number_theory_question
% Stable identity: aops-topic-1371507-q01
\textbf{Six numbers are possible.} Take $4,8,12$ (each $\equiv0\pmod4$) and $1,5,9$ (each $\equiv1\pmod4$). Any four of them include $k$ numbers from $\{1,5,9\}$, where $1\le k\le3$ (there are only three of each kind), and their sum is $\equiv k\not\equiv0\pmod4$.

\textbf{Seven numbers are impossible.} First, among any three integers two have the same parity, so two of them have an even sum. Now take any seven integers.
\begin{itemize}
\item From any three of them pick two with even sum $s_1$; from three of the remaining five pick two with even sum $s_2$; from the remaining three pick two with even sum $s_3$.
\item The three numbers $\dfrac{s_1}{2},\dfrac{s_2}{2},\dfrac{s_3}{2}$ are integers, so two of them, say $\dfrac{s_i}{2}$ and $\dfrac{s_j}{2}$, have the same parity. Then $\dfrac{s_i}{2}+\dfrac{s_j}{2}$ is even, so $s_i+s_j$ is divisible by $4$.
\end{itemize}
The four numbers in these two pairs are four of the seven integers with a sum divisible by $4$. (This is the case $n=4$ of the Erd\H{o}s--Ginzburg--Ziv theorem: among any $2n-1$ integers, some $n$ have a sum divisible by $n$.)

So the largest possible number is $\boxed{6}$.
